\section{Introduction}                   % obligatory
\label{sec:introduction}                 % \secref{sec:introduction}

% ------------------------------------------------------------------------------ host institution
This research project was carried out at ENSTA and more specifically hosted by U2IS (\emph{Unit\'e d'informatique et
d'ing\'enierie des syst\`emes}), the computer science and systems engineering laboratory, whose research covers robotics and artificial intelligence, vision and
perception, and the design of systems with a high level of decision-making autonomy. This project continues a previous internship work on the 3D reconstruction \citep{margaritis2025internship} of a flooded karst conduit
from sonar data and build upon a prior work of \citep{breux2020elevation}. 

% % ------------------------------------------------------------------------------ context and importance
% Flooded karst conduits store the groundwater used by about a quarter of the world's population
% \citep{goldscheider2020karst}, but they also carry dangerous flash floods, as in the city of
% N\^imes \citep{marechal2008nimes}. Managing these aquifers requires knowing the geometry of their
% conduits, which are often too deep, too long or too narrow for human divers. Robots and
% diver-propelled vehicles can go further, and they map the conduits with sonars, since the water is
% often too turbid for cameras \citep{breux2020elevation}. A previous internship built a 3D point
% cloud of the Fontaine de N\^imes conduit from the data of the NavScoot~2
% \citep{margaritis2025internship}, but it used only one of its two sonars: the vertical profiling
% sonar, which measures one cross-section of the conduit at a time. The second sonar, a horizontal
% scanning sonar with a wide beam, sees the walls along the direction of travel and covers much more
% of the conduit at each rotation, but its measurements were left unused because the height of each
% of its echoes inside the $25^\circ$ beam is unknown.

% % ------------------------------------------------------------------------------ objectives
% The aim of this project is to exploit the horizontal sonar. Its objectives are:
% \begin{enumerate}                        % project objectives
%   \item to estimate the elevation angle of every echo of the horizontal sonar, together with its
%         uncertainty, so that each echo becomes a 3D point;
%   \item to do so by implementing the scan-merging method of \citet{breux2020elevation}, which uses
%         the accurate vertical sonar to model the surrounding wall and then finds the elevation at
%         which each horizontal echo lies on that wall;
%   \item to merge the points of both sonars into a single, denser point cloud of the conduit, and to
%         evaluate the result on the data recorded in the Fontaine de N\^imes.
% \end{enumerate}

% % ------------------------------------------------------------------------------ work done
% To reach these objectives, the complete method was implemented in Python and applied to the dive
% of the NavScoot~2, scan by scan. For each complete scan, the trajectory is parameterized by its
% curvilinear abscissa, an elliptic cylinder is fitted to the vertical sonar points as a prior shape
% of the conduit, and a Gaussian process learns how the real wall deviates from this cylinder. Every
% horizontal echo is then tested over the whole vertical aperture of the beam: the most likely
% elevation angles are those whose 3D points fall on the estimated wall, and each of them receives a
% scaled beta distribution describing its uncertainty. Along the way, the pseudocode of the original
% paper was checked line by line and a few errors were corrected. Beyond a denser map for
% hydrogeologists, recovering the 3D information of the wide-beam sonar is a step towards localizing
% an underwater vehicle in a karst by matching successive sonar scans, where the wide beam increases
% the overlap between scans \citep{breux2020elevation}.

% ------------------------------------------------------------------------------ outline
The report is organized as follows. \secref{sec:background} presents the scientific context, the
acquisition platform and the existing processing chain. \secref{sec:definitions} fixes the notation
and models, \secref{sec:msis_merging} describes the scan-merging method and its pseudocode, and
\secref{sec:results} presents the merged point clouds obtained on the Fontaine de N\^imes data.
\secref{sec:conclusions} summarizes the work and proposes directions for future work.